\vspace*{\fill}
\begin{glyph}{Shellfish (貝)}{shellfish}\label{glyph:shellfish}
Shellfish; Seashell.
\end{glyph}

\begin{comps}
\comprtwocone & 目 + 八 \\\hline
\comprthreecone & Eye (\ref{glyph:eye}) + Eight (\ref{glyph:eight})\\\hline
\end{comps}

\begin{remglyph}
\remcom{Cyclops} roasts his \hooglig{shellfish} over a \remcom{volcano}.\\\hline
\end{remglyph}

\begin{prons}
\pronsrtwocone & --- & \textbf{かい (kai)}\\\hline
\pronsrthreecone & バイ (BAI) & ---\\\hline
\end{prons}

\begin{remprons}
\comkun{Kindly} \ucomon{buy} the \hooglig{shellfish}.\\\hline
\end{remprons}

\begin{somme}
貝 & \textit{shellfish} & かい (kai)\\\hline
生貝 & life + shellfish = \textit{raw shellfish} & なま{\middel}がい (nama{\middel}gai)\\\hline
赤貝 & red + shellfish = \textit{ark shell} & あか{\middel}がい (aka{\middel}gai)\\\hline
\end{somme}

\begin{sentence}
\underline{この} \underline{レストラン}の\ruby{\underline{赤貝}}{あかがい}\ruby{は}{わ}\ruby{\underline{有名}}{ユウメイ}{\;}\underline{でぅ}。\\\hline
kono resutoran no aka{\middel}gai wa Y\={U}{\middel}MEI desu.\\\hline
(this) (restaurant) (ark shell) (famous) (is)\\\hline
=\textit{The ark shell in this restaurant is famous.}\\\hline
\end{sentence}
\end{multicols}
\vspace*{-\baselineskip}
\vhrulefill{5pt}
%%--------------------------------------------------------------------
%%--------------------------------------------------------------------
\begin{multicols}{2}
\begin{glyph}{Six (六)}{six}\label{glyph:six}
Six.
\end{glyph}

\begin{comps}
\comprtwocone &  亠 + 八\\\hline
\comprthreecone & Head/Above + Eight (\ref{glyph:eight}) \\\hline
\end{comps}

\begin{remglyph}
The \hooglig{snail} protected itself \remcom{above} the \remcom{volcano}.\\\hline
\end{remglyph}

\begin{prons}
\pronsrtwocone & \textbf{ロク (ROKU)} & ---\\\hline
\pronsrthreecone & リク (RIKU) & むい、むっ、む (mui, mu-, mu)\\\hline
\end{prons}

\begin{remprons}
The \hooglig{6} \ucomkun{mooing} \ucomkun{muslims} were \comon{locked} up to prevent any further \ucomon{licking}.\\\hline
\end{remprons}

\begin{somme}
六 & \textit{six} & ロク (ROKU)\\\hline
六人 & six + person = \textit{six people} & ロク{\middel}ニン (ROKU{\middel}NIN)\\\hline
六月 & six + moon (month) = \textit{June} & ロク{\middel}ガツ (ROKU{\middel}GATSU)\\\hline
六円 & six + circle (yen) = \textit{six yen} & ロク{\middel}エン (ROKU{\middel}EN)\\\hline
六時 & six + time = \textit{six o'clock} & ロク{\middel}ジ (ROKU{\middel}JI)\\\hline
六日 & six + sun (day) = \textit{the sixth day of the month} & むい{\middel}か (mui{\middel}ka)\\\hline
六つ & \textit{six (general counter)} & むっ{\middel}つ (mut{\middel}tsu)\\\hline
\end{somme}

\begin{sentence}
\ruby{\underline{六人}}{ロクニン}\ruby{は}{わ}\ruby{\underline{日本}}{ニホン}\ruby{へ}{え}\ruby{\underline{行}}{い}\underline{きました}。\\\hline
ROKU{\middel}NIN wa NI{\middel}HON e i{\middel}kimashita.\\\hline
(six people) (Japan) (went)\\\hline
=\textit{The six people went to Japan.}\\\hline
\end{sentence}
\end{multicols}
\vspace*{-\baselineskip}
\vhrulefill{5pt}
\vspace*{\fill}
%%--------------------------------------------------------------------
%%--------------------------------------------------------------------
\pagebreak
\vspace*{\fill}
\begin{multicols}{2}
\begin{glyph}{King (王)}{king}\label{glyph:king}
King; royalty.
\end{glyph}

\begin{comps}
\comprtwocone & 王\\\hline
\comprthreecone & King (\ref{glyph:king}) \\\hline
\end{comps}

\begin{remglyph}
---\\\hline
\end{remglyph}

\begin{prons}
\pronsrtwocone & \textbf{オウ (\={O})} & ---\\\hline
\pronsrthreecone & --- & ---\\\hline
\end{prons}

\begin{remprons}
\comon{Awesome} \hooglig{king}.\\\hline
\end{remprons}

\begin{somme}
女王 & woman + king = \textit{queen} & ジョ{\middel}オウ (JO{\middel}\={O})\\\hline
王国 & king + country = \textit{kingdom} & オウ{\middel}コク (\={O}{\middel}KOKU)\\\hline
王朝 & king + morning = \textit{dynasty} & オウ{\middel}チョウ (\={O}{\middel}CH\={O})\\\hline
王子 & king + child = \textit{prince} & オウ{\middel}ジ (\={O}{\middel}JI)\\\hline
海王星 & sea + king + star = \textit{Neptune} & カイ{\middel}オウ{\middel}セイ (KAI{\middel}\={O}{\middel}SEI)\\\hline
\end{somme}

\begin{sentence}
\ruby{\underline{女王}}{ジョオウ}が\ruby{\underline{日本}}{日本}\ruby{へ}{え}\ruby{\underline{行}}{い}\underline{きました}。\\\hline
JO{\middel}\={O} ga NI{\middel}HON e i{\middel}kimashita.\\\hline
(Queen) (Japan) (went)\\\hline
=\textit{The Queen went to Japan.}\\\hline
\end{sentence}
\end{multicols}
\vspace*{-\baselineskip}
\vhrulefill{5pt}
%%--------------------------------------------------------------------
%%--------------------------------------------------------------------
\begin{multicols}{2}
\begin{glyph}{Jewel (玉)}{jewel}\label{glyph:jewel}
Jewel.  Also incorporates the sense of ball-like objects.
\end{glyph}

\begin{comps}
\comprtwocone & 王 + 丶\\\hline
\comprthreecone & King (\ref{glyph:king}) + Dot \\\hline
\end{comps}

\begin{remglyph}
A \remcom{king} carelessly treats his \hooglig{jewels} as if they are random \remcom{dots}.\\\hline
\end{remglyph}

\begin{prons}
\pronsrtwocone & --- & \textbf{たま (tama)}\\\hline
\pronsrthreecone & ギョク (GYOKU) & ---\\\hline
\end{prons}

\begin{remprons}
\comkun{Tamales} is a \hooglig{jewel} of a dish.\\\hline
\end{remprons}

\begin{somme}
玉 & \textit{jewel} & たま (tama)\\\hline
玉ねぎ & \textit{onion} & たま{\middel}ねぎ (tama{\middel}negi)\\\hline
目玉 & eye + jewel = \textit{eyeball} & め{\middel}だま (me{\middel}dama)\\\hline
水玉 & water + jewel = \textit{water droplet} & みず{\middel}たま (mizu{\middel}tama)\\\hline
火玉 & fire + jewel = \textit{fireball} & ひ{\middel}だま (hi{\middel}dama)\\\hline
雪玉 & snow + jewel = \textit{snowball} & ゆき{\middel}だま (yuki{\middel}dama)\\\hline
\end{somme}

\begin{sentence}
\ruby{\underline{玉}}{たま}\underline{ねぎ}\ruby{を}{お}\ruby{\underline{切}}{き}\underline{って} \ruby{\underline{下}}{くだ}\underline{さい}。\\\hline
tama{\middel}negi o ki{\middel}tte kuda{\middel}sai.\\\hline
(onion(s)) (cut) (please).\\\hline
=\textit{Please cut the onion(s)}\\\hline
\end{sentence}
\end{multicols}
\vspace*{-\baselineskip}
\vhrulefill{5pt}
\vspace*{\fill}
%%--------------------------------------------------------------------
%%--------------------------------------------------------------------
\pagebreak
\vspace*{\fill}
\begin{multicols}{2}
\begin{glyph}{Heart (心)}{heart}\label{glyph:heart}
Heart.\\\hline
The sense of feelings.\\\hline
The concept of something's core.
\end{glyph}

\begin{comps}
\comprtwocone & 丶 + 亅 + 丶 + 丶\\\hline
\comprthreecone & Dot + Hook + Dot + Dot \\\hline
\end{comps}

\begin{remglyph}
You can win a \hooglig{heart} by \remcom{doting} or by a \remcom{hook} using brute force.\\\hline
\end{remglyph}

\begin{prons}
\pronsrtwocone & \textbf{シン (SHIN)} & \textbf{こころ (kokoro)}\\\hline
\pronsrthreecone & --- & ---\\\hline
\rye{3}{The kun-yomi for this kanji is always voiced when not in the first position.}
\end{prons}

\begin{remprons}
Her \hooglig{heart} had a pink, romantic \comon{sheen} to it when she gave me the \comkun{Cocoa Roast}.\\\hline
\end{remprons}

\begin{somme}
心 & \textit{heart} & こころ (kokoro)\\\hline
女心 & woman + heart = \textit{a woman's heart} & おんな{\middel}ごころ (onna{\middel}gokoro)\\\hline
中心 & middle + heart = \textit{center} & チュウ{\middel}シン (CH\={U}{\middel}SHIN)\\\hline
下心 & lower + heart = \textit{ulterior motive} & した{\middel}ごころ (shita{\middel}gokoro)\\\hline
安心 & ease + heart = \textit{peace of mind} & アン{\middel}シン (AN{\middel}SHIN)\\\hline
愛国心 & love + country + heart = \textit{patriotism} & アイ{\middel}コク{\middel}シン (AI{\middel}KOKU{\middel}SHIN)\\\hline
\end{somme}

\begin{sentence}
\ruby{\underline{女王}}{ジョオウ}{\;}\underline{によると} \ruby{\underline{愛国心}}{アイコクシン}\ruby{は}{わ}\ruby{\underline{大切}}{タイセツ}だそうです。\\\hline
JO{\middel}\={O} ni yoru to AI{\middel}KOKU{\middel}SHIN wa TAI{\middel}SETSU da s\={o} desu.\\\hline
(queen) (according to) (patriotism) (important)\\\hline
=\textit{According to the queen, patriotism is important.}\\\hline
\end{sentence}
\end{multicols}
\vspace*{-\baselineskip}
\vhrulefill{5pt}
%%--------------------------------------------------------------------
%%--------------------------------------------------------------------
\begin{multicols}{2}
\begin{glyph}{Country (国)}{country}\label{glyph:country}
Country; Nation; State.
\end{glyph}

\begin{comps}
\comprtwocone & 囗 + 玉\\\hline
\comprthreecone & Prison/Enclosure + Jewel (\ref{glyph:jewel}) \\\hline
\end{comps}

\begin{remglyph}
Within these \remcom{borders} is a \remcom{jewel} of a \hooglig{country}.\\\hline
\end{remglyph}

\begin{prons}
\pronsrtwocone & \textbf{コク (KOKU)} & \textbf{くに (kuni)}\\\hline
\pronsrthreecone & --- & ---\\\hline
\end{prons}

\begin{remprons}
That \hooglig{country's} \comkun{cuneiform} \comon{cocked} my interest.\\\hline
\end{remprons}

\begin{somme}
国 & \textit{country} & コク (KOKU)\\\hline
王国 & king + country = \textit{Kingdom} & オウ{\middel}コク (\={O}{\middel}KOKU)\\\hline
全国 & complete + country = \textit{the whole country/nationwide} & ゼン{\middel}コク (ZEN{\middel}KOKU)\\\hline
入国 & enter + country = \textit{to enter a country} & ニュウ{\middel}コク (NY\={U}{\middel}KOKU)\\\hline
国内 & country + inside = \textit{domestic} & コク{\middel}ナイ (KOKU{\middel}NAI)\\\hline
外国人 & outside + country + person = \textit{foreigner} & ガイ{\middel}コク{\middel}ジン (GAI{\middel}KOKU{\middel}JIN)\\\hline
愛国心 & love + country + heart = \textit{patriotism} & アイ{\middel}コク{\middel}シン (AI{\middel}KOKU{\middel}SHIN)\\\hline
\end{somme}

\begin{sentence}
\underline{あの} \ruby{\underline{国}}{くに}には\ruby{\underline{外国人}}{ガイコクジン}が\ruby{\underline{多}}{おお}\underline{い}。\\\hline
ano kuni ni wa GAI{\middel}KOKU{\middel}JIN ga \={o}{\middel}i.\\\hline
(that) (country) (foreigners) (many)\\\hline
=\textit{There are many foreigners in that country.}\\\hline
\end{sentence}
\end{multicols}
\vspace*{-\baselineskip}
\vhrulefill{5pt}
\vspace*{\fill}
%%--------------------------------------------------------------------
%%--------------------------------------------------------------------
\vspace*{\fill}
\pagebreak
\begin{multicols}{2}
\begin{glyph}{Complete (全)}{complete}\label{glyph:complete}
Complete; Wholeness; Completion.
\end{glyph}

\begin{comps}
\comprtwocone & 𠆢 + 王 \\\hline
\comprthreecone & Person (\ref{glyph:person}) + King (\ref{glyph:king}) \\\hline
\end{comps}

\begin{remglyph}
A \remcom{king} is a \hooglig{complete} \remcom{person}.\\\hline
\end{remglyph}

\begin{prons}
\pronsrtwocone & \textbf{ゼン (ZEN)} & --- \\\hline
\pronsrthreecone & --- & まった、まっと (matta, matto) \\\hline
\end{prons}

\begin{remprons}
I had \hooglig{complete} \comon{Zen} until the \ucomkun{mutt} attacked me with a \ucomkun{mattock}.\\\hline
\end{remprons}

\begin{somme}
全国 & complete + country = \textit{the whole country/nationwide} & ゼン{\middel}コク (ZEN{\middel}KOKU)\\\hline
安全 & ease + complete = \textit{safety} & アン{\middel}ゼン (AN{\middel}ZEN)\\\hline
全校 & complete + school = \textit{the whole school/all schools} & ゼン{\middel}コウ (ZEN{\middel}K\={O})\\\hline
全力 & complete + strength = \textit{all one's power} & ゼン{\middel}リョク (ZEN{\middel}RYOKU)\\\hline
全体 & complete + body = \textit{the whole} & ゼン{\middel}タイ (ZEN{\middel}TAI)\\\hline
全部 & complete + part = \textit{all} & ゼン{\middel}ブ (ZEN{\middel}BU)\\\hline
\end{somme}

\begin{sentence}
\ruby{\underline{田中}}{たなか}\underline{さん}\ruby{は}{わ}\ruby{\underline{全国}}{ゼンコク}\ruby{を}{お}\ruby{\underline{回}}{まわ}\underline{る}。\\\hline
Ta{\middel}naka-san wa ZEN{\middel}KOKU o mawa{\middel}ru.\\\hline
(Tanaka-san) (whole country) (go around)\\\hline
=\textit{Tanaka-san is going around the whole country.}\\\hline
\end{sentence}
\end{multicols}
\vspace*{-\baselineskip}
\vhrulefill{5pt}
%%--------------------------------------------------------------------
%%--------------------------------------------------------------------
\begin{multicols}{2}
\begin{glyph}{Ten (十)}{ten}\label{glyph:ten}
Ten.\\\hline
The first two examples of Entry 176 show another use for this character, an interesting application based entirely on its shape.
\end{glyph}

\begin{comps}
\comprtwocone & 十 \\\hline
\comprthreecone & Ten (\ref{glyph:ten}) \\\hline
\end{comps}

\begin{remglyph}
\remcom{Scarecrow}.\\\hline
\end{remglyph}

\begin{prons}
\pronsrtwocone & \textbf{ジュウ (J\={U})} & --- \\\hline
\pronsrthreecone & ジツ (JITSU) & とお、と (t\={o}, to) \\\hline
\end{prons}

\begin{remprons}
As the Brave \hooglig{10} were \ucomkun{tormented} by the \ucomkun{torrents} of rain, they had to get their creative \comon{juices} flowing to keep on practising \ucomon{ju-jitsu}.\\\hline
\end{remprons}

\begin{somme}
十 & \textit{ten} & ジュウ (J\={U})\\\hline
十月 & ten + moon (month) = \textit{October} & ジュウ{\middel}ガツ (J\={U}{\middel}GATSU)\\\hline
十一月 & ten + one + moon (month) = \textit{November} & ジュウ{\middel}イチ{\middel}ガツ (J\={U}{\middel}ICHI{\middel}GATSU)\\\hline
十二月 & ten + two + moon (month) = \textit{December} & ジュウ{\middel}ニ{\middel}ガツ (J\={U}{\middel}NI{\middel}GATSU)\\\hline
十六 & ten + six = \textit{sixteen} & ジュウ{\middel}ロク (J\={U}{\middel}ROKU)\\\hline
十八 & ten + eight = \textit{eighteen} & ジュウ{\middel}ハチ (J\={U}{\middel}HACHI)\\\hline
十日 & ten + sun (day) = \textit{the tenth day of the month} & とお{\middel}か (t\={o}{\middel}ka)\\\hline
\end{somme}

\begin{irrr}
二十日 & two + ten + sun (day) = \textit{the twentieth day of the month} & はつか (hatsuka)\\\hline
二十歳 & two + ten + annual = \textit{twenty years old} & はたち (hatachi)\\\hline
\end{irrr}

\begin{sentence}
\ruby{\underline{十月}}{ジュウガツ}に\ruby{は}{わ}\underline{あの} \ruby{\underline{木}}{き}\ruby{は}{わ}\ruby{\underline{美}}{うつく}\underline{しく} \underline{なります}。\\\hline
J\={U}{\middel}GATSU ni wa ano ki wa utsuku{\middel}shiku narimasu.\\\hline
(October) (that) (tree) (beautiful) (becomes)\\\hline
=\textit{That tree becomes beautiful in October.}\\\hline
\end{sentence}
\end{multicols}
\vspace*{-\baselineskip}
\vhrulefill{5pt}
\vspace*{\fill}
%%--------------------------------------------------------------------
%%--------------------------------------------------------------------
\pagebreak
\vspace*{\fill}
\begin{multicols}{2}
\begin{glyph}{Early (早)}{early}\label{glyph:early}
Early; Fast.
\end{glyph}

\begin{comps}
\comprtwocone & 日 + 十 \\\hline
\comprthreecone & Sun (\ref{glyph:sun}) + Ten (\ref{glyph:ten}) \\\hline
\end{comps}

\begin{remglyph}
\hooglig{Early} in the \remcom{morning-sun} the \remcom{scarecrow} gets into shape.\\\hline
\end{remglyph}

\begin{prons}
\pronsrtwocone & \textbf{ソウ (S\={O})} & \textbf{はや (haya)} \\\hline
\pronsrthreecone & サツ (SATSU) & --- \\\hline
\end{prons}

\begin{remprons}
The \ucomon{sutler soup} eased the \comon{sorely} displeasing \comkun{hay allergy} I had \hooglig{earlier}.\\\hline
\end{remprons}

\begin{somme}
早い & \textit{early; fast} & はや{\middel}い (haya{\middel}i)\\\hline
早まる & \textit{to be in a hurry} & はや{\middel}まる (haya{\middel}maru)\\\hline
早める & \textit{to hurry something up} & はや{\middel}める (haya{\middel}meru)\\\hline
早口 & early + mouth = \textit{rapid speaking} & はや{\middel}くち (haya{\middel}kuchi)\\\hline
早春 & early + spring = \textit{early spring} & ソウ{\middel}シュン (S\={O}{\middel}SHUN)\\\hline
早朝 & early + morning = \textit{early morning} & ソウ{\middel}チョウ (S\={O}{\middel}CH\={O})\\\hline
\end{somme}

\begin{sentence}
\underline{どうして} \underline{そんな}に\ruby{\underline{早}}{はや}\underline{く} \ruby{\underline{来}}{き}\underline{た}のですか。\\\hline
d\={o}shite sonna ni haya{\middel}ku ki{\middel}ta no desu ka.\\\hline
(why) (so) (early) (come)\\\hline
=\textit{Why did you come so early?}\\\hline
\end{sentence}
\end{multicols}
\vspace*{-\baselineskip}
\vhrulefill{5pt}
%%--------------------------------------------------------------------
%%--------------------------------------------------------------------
\begin{multicols}{2}
\begin{glyph}{Upper (上)}{upper}\label{glyph:upper}
Upper; On; Over.\\\hline
It can refer to anything from goods of high quality to superiors at work.
\end{glyph}

\begin{comps}
\comprtwocone & 卜 + 一 \\\hline
\comprthreecone & Divination + One (\ref{glyph:one}) \\\hline
\end{comps}

\begin{remglyph}
\remcom{One} \hooglig{upper} \remcom{divination}.\\\hline
\end{remglyph}

\begin{prons}
\pronsrtwocone & \textbf{ジョウ (J\={O})} & \textbf{かみ、うえ、のぼ、あ、うわ (kami, ue, nobo, a, uwa)} \\\hline
\pronsrthreecone & ショウ (SH\={O}) & --- \\\hline
\rye{3}{J\={O} is by far the most frequently used pronounciation.}
\end{prons}

\begin{remprons}
\hooglig{Upper} \comon{jaw}.\\\hline
\end{remprons}

\begin{somme}
上 & \textit{upper} & かみ (kami)\\\hline
上 & \textit{over; on} & うえ (ue)\\\hline
上る (intr) & \textit{to climb} & のぼ{\middel}る (nobo{\middel}ru)\\\hline
上がる (intr) & \textit{to rise} & あ{\middel}がる (a{\middel}garu)\\\hline
上げる (tr/intr) & \textit{to lift} & あ{\middel}げる (a{\middel}geru)\\\hline
上手 & upper + hand = \textit{skillful} & じょう{\middel}ズ (J\={O}{\middel}ZU)\\\hline
\rye{3}{The verb あ{\middel}げる is almost always used in a transitive sense (that is, it lifts some objects in either a physical or symbolic way), and is best thought of as being paired with the intransitive あ{\middel}がる.}
\end{somme}

\begin{sentence}
\ruby{\underline{肉}}{ニク}\ruby{は}{わ}\underline{テーベル}の\ruby{\underline{上}}{うえ}に\underline{あります}。\\\hline
NIKU wa T\={E}BERU no ue ni arimasu.\\\hline
(meat) (table) (on) (is)\\\hline
=\textit{The meat is on the table.}\\\hline
\end{sentence}
\end{multicols}
\vspace*{-\baselineskip}
\vhrulefill{5pt}
\vspace*{\fill}
%%--------------------------------------------------------------------
%%--------------------------------------------------------------------
\pagebreak
\begin{multicols}{2}
\begin{glyph}{Lower (下)}{lower}\label{glyph:lower}
Lower; Under; Below.
\end{glyph}

\begin{comps}
\comprtwocone & 一 + 卜 \\\hline
\comprthreecone & One (\ref{glyph:one}) + Divination \\\hline
\end{comps}

\begin{remglyph}
\remcom{One} \hooglig{lower} \remcom{divination}.\\\hline
\end{remglyph}

\begin{prons}
\pronsrtwocone & \textbf{カ、ゲ (KA, GE)} & \textbf{した、しま、お、さ、くだ (shita, shimo, o, sa, kuda)}\\\hline
\pronsrthreecone & --- & もと (moto) \\\hline
\rye{3}{お、さ、くだ above are all verb stems, and will thus be accompanied by hiragana hinting at the correct pronounciation.\\\hline Regarding compounds, カ and ゲ will be encountered far more than the others, although there is no obvious pattern as to when each of these is used.}
\end{prons}

\begin{remprons}
---\\\hline
\end{remprons}

\begin{somme}
下 & \textit{low; below; under} & した (shita)\\\hline
下りる (intr) & \textit{to come down} & お{\middel}りる (o{\middel}riru)\\\hline
下ろす (tr) & \textit{to take down} & お{\middel}ろす (o{\middel}rosu)\\\hline
下がる (intr) & \textit{to hang down (on one's own)} & さ{\middel}がる (sa{\middel}garu)\\\hline
下げる (tr) & \textit{to lower (something)} & さ{\middel}げる (sa{\middel}geru)\\\hline
下さる & \textit{to oblige} & くだ{\middel}さる (kuda{\middel}saru)\\\hline
上下 & upper + lower = \textit{high and low} & ジョウ{\middel}ゲ (J\={O}{\middel}GE)\\\hline
天下 & heaven + lower = \textit{the whole land} & テン{\middel}カ (TEN{\middel}KA)\\\hline
\end{somme}

\begin{irrr}
下手 & lower + hand = \textit{to be poor at something} & へ{\middel}た (he{\middel}ta)\\\hline
\end{irrr}

\begin{sentence}
\ruby{\underline{木}}{き}{\;}\underline{から} \ruby{\underline{下}}{お}\underline{りて} \ruby{\underline{下}}{くだ}\underline{さい}。\\\hline
ki kara o{\middel}rite kuda{\middel}sai.\\\hline
(tree) (from) (come down) (please)\\\hline
=\textit{Please come down from the tree.}\\\hline
\end{sentence}
\end{multicols}
\vspace*{-\baselineskip}
\vhrulefill{5pt}
%%--------------------------------------------------------------------
%%--------------------------------------------------------------------
\begin{multicols}{2}
\begin{glyph}{Rice (米)}{rice}\label{glyph:rice}
(Uncooked) rice.\\\hline
This character is also used to symbolize the America's.
\end{glyph}

\begin{comps}
\comprtwocone & 米 (\ref{glyph:rice}) \\\hline
\comprthreecone & Bunny Ears + Tree (\ref{glyph:tree}) \\\hline
\end{comps}

\begin{remglyph}
---\\\hline
\end{remglyph}

\begin{prons}
\pronsrtwocone & \textbf{マイ (MAI)} & \textbf{こめ (kome)}\\\hline
\pronsrthreecone & ベイ (BEI) & --- \\\hline
\rye{3}{ベイ is used as the reading in compounds relating to the Americas.\\\hline
マイ is the primary choice when it refers to rice.}
\end{prons}

\begin{remprons}
We'll make our \comkun{comeback} by \comon{basically} \comon{migrating} to regions with plentiful \hooglig{rice}.\\\hline
\end{remprons}

\begin{somme}
米 & \textit{(uncooked) rice} & こめ (kome)\\\hline
白米 & white + rice = \textit{polished rice} & ハク{\middel}マイ (HAKU{\middel}MAI)\\\hline
米国 & rice + country = \textit{The United States} & ベイ{\middel}コク (BEI{\middel}KOKU)\\\hline
日米 & sun + rice = \textit{Japan-U.S} & ニチ{\middel}ベイ (NICHI{\middel}BEI)\\\hline
中米 & middle + rice = \textit{Central America} & チュウ{\middel}ベイ (CH\={U}{\middel}BEI)\\\hline
北米 & North + rice = \textit{North America} & ホク{\middel}ベイ (HOKU{\middel}BEI)\\\hline
南米 & South + rice = \textit{South America} & ナン{\middel}ベイ (NAN{\middel}BEI)\\\hline
\end{somme}

\begin{sentence}
\ruby{\underline{米国}}{ベイコク}の\ruby{\underline{牛肉}}{ギュウニク}\ruby{は}{わ}\ruby{\underline{高}}{たか}\underline{い}。\\\hline
BEI{\middel}KOKU no GY\={U}{\middel}NIKU wa taka{\middel}i.\\\hline
(United States) (beef) (expensive)\\\hline
=\textit{U.S. beef is expensive.}\\\hline
\end{sentence}
\end{multicols}
\vspace*{-\baselineskip}
\vhrulefill{5pt}
%%--------------------------------------------------------------------
%%--------------------------------------------------------------------
\pagebreak
\vspace*{\fill}
\begin{multicols}{2}
\begin{glyph}{Self (自)}{self}\label{glyph:self}
Self.
\end{glyph}

\begin{comps}
\comprtwocone & 丶 + 目 \\\hline
\comprthreecone & Dot + Eye (\ref{glyph:eye}) \\\hline
\end{comps}

\begin{remglyph}
\remcom{Cyclops} loses \hooglig{self}-control when he sees a \remcom{dot}.\\\hline
\end{remglyph}

\begin{prons}
\pronsrtwocone & \textbf{ジ (JI)} & ---\\\hline
\pronsrthreecone & シ (SHI) & みずか (mizuka) \\\hline
\rye{3}{シ only appears in the compound 自然: self + nature = \textit{Nature}, シ{\middel}ゼン.}
\end{prons}

\begin{remprons}
I told \hooglig{myself}: ``Make \ucomkun{me zoom completely} on the \comon{jigged} \ucomon{shields}''.\\\hline
\end{remprons}

\begin{somme}
自国 & self + country = \textit{one's homeland} & ジ{\middel}コク (JI{\middel}KOKU)\\\hline
自立 & self + stand = \textit{independence} & ジ{\middel}リツ (JI{\middel}RITSU)\\\hline
自分 & self + part = \textit{oneself} & ジ{\middel}ブン (JI{\middel}BUN)\\\hline
自体 & self + body = \textit{in itself} & ジ{\middel}タイ (JI{\middel}TAI)\\\hline
自然 & self + body = \textit{nature} & シ{\middel}ゼン (SHI{\middel}ZEN)\\\hline
\end{somme}

\begin{sentence}
\ruby{\underline{自分}}{ジブン}の\ruby{\underline{車}}{くるま}が\ruby{\underline{買}}{か}\underline{いたい}。\\\hline
JI{\middel}BUN no kuruma ga ka{\middel}itai.\\\hline
(oneself) (car) (want to buy)\\\hline
=\textit{I want to buy my own car.}\\\hline
\end{sentence}
\end{multicols}
\vspace*{-\baselineskip}
\vhrulefill{5pt}
%%--------------------------------------------------------------------
%%--------------------------------------------------------------------
\begin{multicols}{2}
\begin{glyph}{Inside (内)}{inside}\label{glyph:inside}
Inside; Interior; Within.\\\hline
This kanji can also convey the idae of ``home''.
\end{glyph}

\begin{comps}
\comprtwocone & 冂 + 人 \\\hline
\comprthreecone & Wilderness + Person (\ref{glyph:person}) \\\hline
\end{comps}

\begin{remglyph}
The \hooglig{inside} of a \remcom{person} is a \remcom{wilderness}.\\\hline
\end{remglyph}

\begin{prons}
\pronsrtwocone & \textbf{ナイ (NAI)} & \textbf{うち (uchi)}\\\hline
\pronsrthreecone & ダイ (DAI) & --- \\\hline
\end{prons}

\begin{remprons}
\hooglig{In} the \comkun{Uchica} community are \comon{nicely} \ucomon{dyed} fabrics.\\\hline
\end{remprons}

\begin{somme}
内 & \textit{inside} & うち (うち)\\\hline
国内 & country + inside = \textit{domestic} & コク{\middel}ナイ (KOKU{\middel}NAI)\\\hline
市内 & city + inside = \textit{within the city} & シ{\middel}ナイ (SHI{\middel}NAI)\\\hline
内海 & inside + sea = \textit{inland sea} & ナイ{\middel}カイ (NAI{\middel}KAI)\\\hline
車内 & car + inside = \textit{inside the car} & シャ{\middel}ナイ (SHA{\middel}NAI)\\\hline
町内 & town + inside = \textit{in the town} & チョウ{\middel}ナイ (CH\={O}{\middel}NAI)\\\hline
\end{somme}

\begin{sentence}
\underline{その} \ruby{\underline{人}}{ひと}\ruby{は}{わ}\ruby{\underline{国内}}{コクナイ}で\ruby{\underline{有名}}{ユウメイ}です。\\\hline
sono hito wa KOKU{\middel}NAI de Y\={U}{\middel}MEI desu.\\\hline
(that) (person) (domestic) (famous)\\\hline
=\textit{That person is famous domestically.}\\\hline
\end{sentence}
\end{multicols}
\vspace*{-\baselineskip}
\vhrulefill{5pt}
\vspace*{\fill}
%%--------------------------------------------------------------------
%%--------------------------------------------------------------------
\pagebreak
\vspace*{\fill}
\begin{multicols}{2}
\begin{glyph}{Right (右)}{migi}\label{glyph:right}
Right; Right-hand side.
\end{glyph}

\begin{comps}
\comprtwocone & 一 + ノ + 口  \\\hline
\comprthreecone & One (\ref{glyph:one}) + No (katakana)/Slash + Mouth (\ref{glyph:mouth}) \\\hline
\end{comps}

\begin{remglyph}
\hooglig{Rightly} \remcom{slash} that \remcom{one's} \remcom{mouth}.\\\hline
\end{remglyph}

\begin{prons}
\pronsrtwocone & --- & \textbf{みぎ (migi)}\\\hline
\pronsrthreecone & ユウ、ウ (Y\={U}, U) & --- \\\hline
\end{prons}

\begin{remprons}
On \ucomon{Youtube} was an \ucomon{ooky} clip of a guy \comkun{midging} with his \hooglig{right} hand.\\\hline
\end{remprons}

\begin{somme}
右 & \textit{right} & みぎ (migi)\\\hline
右回り & right + rotate = \textit{clockwise} & みぎ　まわ{\middel}り (migi mawa{\middel}ri)\\\hline
右上 & right + upper = \textit{upper right} & みぎ{\middel}うえ (migi{\middel}ue)\\\hline
右下 & right + lower = \textit{lower right} & みぎ{\middel}した (migi{\middel}shita)\\\hline
右手 & right + hand = \textit{right hand} & みぎ{\middel}て (migi{\middel}te)\\\hline
右左 & right + leg = \textit{right leg} & みぎ{\middel}あし (migi{\middel}ashi)\\\hline
\end{somme}

\begin{sentence}
\ruby{\underline{山口}}{やまぐち}\underline{さん}\ruby{は}{わ}\ruby{\underline{右}}{みぎ}の\ruby{\underline{目}}{め}が\underline{かゆい}。\\\hline
Yama{\middel}guchi-san wa migi no me ga kayui.\\\hline
(Yamaguchi-san) (right) (eye) (itchy)\\\hline
=\textit{Yamaguchi-san's right eye is itchy.}\\\hline
\end{sentence}
\end{multicols}
\vspace*{-\baselineskip}
\vhrulefill{5pt}
%%--------------------------------------------------------------------
%%--------------------------------------------------------------------
\begin{multicols}{2}
\begin{glyph}{Have (有)}{have}\label{glyph:have}
Have; Possess.
\end{glyph}

\begin{comps}
\comprtwocone & 一 + ノ + 月  \\\hline
\comprthreecone & One (\ref{glyph:one}) + No (katakana)/Slash + Moon (\ref{glyph:moon}) \\\hline
\end{comps}

\begin{remglyph}
\hooglig{Have} \remcom{one} \remcom{slash} at the \remcom{moon}.\\\hline
\end{remglyph}

\begin{prons}
\pronsrtwocone & \textbf{ユウ (Y\={U})} & \textbf{あ (a)}\\\hline
\pronsrthreecone & ウ (U) & --- \\\hline
\end{prons}

\begin{remprons}
Many people on \comon{Youtube} \hooglig{have} \ucomon{ooky} \comkun{updates}.\\\hline
\end{remprons}

\begin{somme}
有る & \textit{to have} & あ{\middel}る (a{\middel}ru)\\\hline
有力 & have + strength = \textit{powerful} & ユウ{\middel}リョク (Y\={U}{\middel}RYOKU)\\\hline
有名 & have + name = \textit{famous} & ユウ{\middel}メイ (Y\={U}{\middel}MEI)\\\hline
私有 & private + have = \textit{privately owned} & シ{\middel}ユウ (SHI{\middel}Y\={U})\\\hline
公有 & public + have = \textit{publicly owned} & コウ{\middel}ユウ (K\={O}{\middel}Y\={U})\\\hline
国有化 & country + have + change = \textit{nationalization} & コク{\middel}ユウ{\middel}カ (KOKU{\middel}Y\={U}{\middel}KA)\\\hline
\end{somme}

\begin{sentence}
\underline{エジプト}の\underline{ピラミッド}\ruby{は}{わ}\ruby{\underline{有名}}{なまえ}です。\\\hline
Ejiputo no piramiddo wa Y\={U}{\middel}MEI desu.\\\hline
(Egypt) (pyramids) (famous)\\\hline
=\textit{The Egyptian pyramids are famous.}\\\hline
\end{sentence}
\end{multicols}
\vspace*{-\baselineskip}
\vhrulefill{5pt}
\vspace*{\fill}
%%--------------------------------------------------------------------
%%--------------------------------------------------------------------
\pagebreak
\vspace*{\fill}
\begin{multicols}{2}
\begin{glyph}{Meat (肉)}{meat}\label{glyph:meat}
Meat; Flesh.\\\hline
This kanji also has a sexual connotation through its appearance in words such as ``lust'' and ``sensuality''.
\end{glyph}

\begin{comps}
\comprtwocone & 内 + 人  \\\hline
\comprthreecone & Inside (\ref{glyph:inside}) + Person (\ref{glyph:person}) \\\hline
\end{comps}

\begin{remglyph}
There's \hooglig{meat} \remcom{inside} a \remcom{person}.\\\hline
\end{remglyph}

\begin{prons}
\pronsrtwocone & \textbf{ニク (NIKU)} & ---\\\hline
\pronsrthreecone & --- & --- \\\hline
\end{prons}

\begin{remprons}
The \hooglig{flesh} lusts after \comon{nicotine}.\\\hline
\end{remprons}

\begin{somme}
肉 & \textit{meat} & ニク (NIKU)\\\hline
人肉 & person + meat = \textit{human flesh} & ジン{\middel}ニク (JIN{\middel}NIKU)\\\hline
牛肉 & cow + meat = \textit{beef} & ギュウ{\middel}ニク (GY\={U}{\middel}NIKU)\\\hline
肉親 & meat + parent = \textit{blood relative} & ニク{\middel}シン (NIKU{\middel}SHIN)\\\hline
肉体 & meat + body = \textit{the body} & ニク{\middel}タイ (NIKU{\middel}TAI)\\\hline
\end{somme}

\begin{sentence}
\underline{あの} \ruby{\underline{四人}}{よニン}\ruby{は}{わ}\ruby{\underline{牛肉}}{ギュウニク}が\ruby{\underline{大好}}{ダイす}\underline{き}です。\\\hline
ano yo{\middel}NIN wa GY\={U}{\middel}NIKU ga DAIsu{\middel}ki desu.\\\hline
(those) (four people) (beef) (really like)\\\hline
=\textit{Those four people really like beef.}\\\hline
\end{sentence}
\end{multicols}
\vspace*{-\baselineskip}
\vhrulefill{5pt}
%%--------------------------------------------------------------------
%%--------------------------------------------------------------------
\begin{multicols}{2}
\begin{glyph}{Half (半)}{half}\label{glyph:half}
Half.\\\hline
This character also incorporates a sense of ``partial'', forming many words that English conveys with the prefix ``semi-''.
\end{glyph}

\begin{comps}
\comprtwocone & 丶丶 + 二 + 丨\\\hline
\comprthreecone & Dots + Two (\ref{glyph:two}) + Pole \\\hline
\end{comps}

\begin{remglyph}
A \remcom{pole} can \hooglig{half} a the \remcom{two dots}.\\\hline
\end{remglyph}

\begin{prons}
\pronsrtwocone & \textbf{ハン (HAN)} & ---\\\hline
\pronsrthreecone & --- & なか (naka) \\\hline
\end{prons}

\begin{remprons}
I have a \comon{hunch} he has a \ucomkun{knack} at \hooglig{halving} wood.\\\hline
\end{remprons}

\begin{somme}
上半 & upper + half = \textit{upper half} & ジョウ{\middel}ハン (J\={O}{\middel}HAN)\\\hline
下半 & lower + half = \textit{lower half} & カ{\middel}ハン (KA{\middel}HAN)\\\hline
半日 & half + sun (day) = \textit{half a day} & ハン{\middel}ニチ (HAN{\middel}NICHI)\\\hline
半分 & half + part = \textit{half} & ハン{\middel}ブン (HAN{\middel}BUN)\\\hline
九時半 & nine + time + half = \textit{nine-thirty} & キュウ{\middel}ジ{\middel}ハン (KY\={U}{\middel}JI{\middel}HAN)\\\hline
半円 & half + circle = \textit{semicircle} & ハン{\middel}エン (HAN{\middel}EN)\\\hline
\end{somme}

\begin{sentence}
\ruby{\underline{木口}}{きぐち}\underline{さん}\ruby{は}{わ}\ruby{\underline{半日}}{ハンニチ}{\;}\underline{ぐらい} \ruby{\underline{新聞}}{シンブン}\ruby{を}{お}\ruby{\underline{読}}{よ}\underline{みました}。\\\hline
Ki{\middel}guchi-san wa HAN{\middel}NICHI gurai SHIN{\middel}BUN o yo{\middel}mimashita.\\\hline
(Kiguchi-san) (half a day) (about) (newspaper) (read)\\\hline
=\textit{Kiguchi-san read the newspaper for about half a day.}\\\hline
\end{sentence}
\end{multicols}
\vspace*{-\baselineskip}
\vhrulefill{5pt}
\vspace*{\fill}
%%--------------------------------------------------------------------
%%--------------------------------------------------------------------
\pagebreak
\vspace*{\fill}
\begin{multicols}{2}
\begin{glyph}{Nine (九)}{nine}\label{glyph:nine}
Nine.
\end{glyph}

\begin{comps}
\comprtwocone & 九　\\\hline
\comprthreecone & Nine (\ref{glyph:nine}) \\\hline
\end{comps}

\begin{remglyph}
---\\\hline
\end{remglyph}

\begin{prons}
\pronsrtwocone & \textbf{キュウ、ク (KY\={U}, KU)} & ---\\\hline
\pronsrthreecone & --- & ここの (kokono) \\\hline
\rye{3}{キュウ is the general reading for this kanji.\\\hline
Only likely to encounter ク in the third and fourth examples below.}
\end{prons}

\begin{remprons}
\comon{Cooking} \ucomkun{coconut} at \hooglig{nine o'clock} will \comon{kill you}.\\\hline
\end{remprons}

\begin{somme}
九 & \textit{nine} & キュウ (KY\={U})\\\hline
九十 & nine + ten = \textit{ninety} & キュウ{\middel}ジュウ (KY\={U}{\middel}J\={U})\\\hline
九月 & nine + moon (month) = \textit{September} & ク{\middel}ガツ (KU{\middel}GATSU)\\\hline
九時 & nine + time = \textit{nine o'clock} & キュウ{\middel}ジ (KU{\middel}JI)\\\hline
\end{somme}

\begin{sentence}
\ruby{\underline{中山}}{なかやま}\underline{さん}\ruby{は}{わ}\ruby{\underline{九月}}{クガツ}に\ruby{\underline{日本}}{ニホン}\ruby{へ}{え}\ruby{\underline{行}}{い}\underline{きました}。\\\hline
Naka{\middel}yama-san wa KU{\middel}GATSU ni NI{\middel}HON e i{\middel}kimashita.\\\hline
(Nakayama-san) (September) (Japan) (went)\\\hline
=\textit{Nakayama-san went to Japan in September.}\\\hline
\end{sentence}
\end{multicols}
\vspace*{-\baselineskip}
\vhrulefill{5pt}
%%--------------------------------------------------------------------
%%--------------------------------------------------------------------
\begin{multicols}{2}
\begin{glyph}{Spring (春)}{spring}\label{glyph:spring}
Spring.\\\hline
This kanji also carries a hint of sexuality and amorousness---this is apparent by its presence in the Japanese words for puberty, etc.
\end{glyph}

\begin{comps}
\comprtwocone & 三 + 人 + 日　\\\hline
\comprthreecone & Three (\ref{glyph:three}) + Person (\ref{glyph:person}) + Sun (\ref{glyph:sun}) \\\hline
\end{comps}

\begin{remglyph}
\hooglig{Spring} is the \remcom{third} \remcom{personage} of the \remcom{sun}.\\\hline
\end{remglyph}

\begin{prons}
\pronsrtwocone & \textbf{シュン (SHUN)} & \textbf{はる (haru)}\\\hline
\pronsrthreecone & --- & --- \\\hline
\end{prons}

\begin{remprons}
\comon{Shun} the \hooglig{spring} \comkun{haruspex}.\\\hline
\end{remprons}

\begin{somme}
春 & \textit{spring} & はる (haru)\\\hline
小春 & small + spring = \textit{Indian Summer} & こ{\middel}はる (ko{\middel}haru)\\\hline
青春 & Blue + spring = \textit{youth} & セイ{\middel}シュン (SEI{\middel}SHUN)\\\hline
売春 & sell + spring = \textit{prostitution} & バイ{\middel}シュン (BAI{\middel}SHUN)\\\hline
春分 & spring + part = \textit{vernal equinox} & シュン{\middel}ブン (SHUN{\middel}BUN)\\\hline
来春 & come + spring = \textit{next spring} & ライ{\middel}シュン (RAI{\middel}SHUN)\\\hline
\end{somme}

\begin{sentence}
\ruby{\underline{日本}}{ニホン}の\ruby{\underline{春}}{はる}\ruby{は}{わ}\ruby{\underline{美}}{うつく}\underline{しい}。\\\hline
NI{\middel}HON no haru wa utsuku{\middel}shii.\\\hline
(Japan) (spring) (beautiful)\\\hline
=\textit{Japan's spring is beautiful.}\\\hline
\end{sentence}
\end{multicols}
\vspace*{-\baselineskip}
\vhrulefill{5pt}
\vspace*{\fill}
%%--------------------------------------------------------------------
%%--------------------------------------------------------------------
\pagebreak